\section{Simulating DESI-like Images and Spectra}
\label{sec:data-generate}

\subsection{Galaxy and Observation Model}
\label{ssec:kl-model}

We simulate our galaxies following the DESI-KL concept of \citep{KLDESI}, in which each galaxy is observed with a photometric image and a set of DESI fiber spectra. DESI takes spectra by placing optical fibers of 1.5$\arcsec$ diameter on its targets \citep{DESIOverview}. For KL, we need multiple fibers placed along a disk galaxy's major and minor axes in order to constrain its velocity field.

A similar fiber arrangement is already used by the DESI Peculiar Velocity Survey (DESI-PV), which places offset fibers along the major axes of disk galaxies to measure their rotation velocities \citep{DESIPVTargetSel2023, DESIPVTF2025}. This makes DESI-PV data a natural candidate for a first KL measurement with DESI, and such a pilot measurement is being carried out in Willett et al. (in prep.). The method presented in this work could eventually be applied to such data as well.

\begin{table}
    \centering
    \begin{tabular}{llcl}
        \toprule
        Parameter & Description & Training prior & Units\\
        \tableline
        $g_1$ & Reduced shear & $\mathcal{U}$(-0.1, 0.1) & 1\\
        $g_2$ & Reduced shear & $\mathcal{U}$(-0.1, 0.1) & 1\\
        $\theta_\mathrm{int}$ & Position angle & $\mathcal{U}$(-$\pi$, $\pi$) & rad\\
        $\cos i$ & Inclination & $\mathcal{U}$(0, 1) & 1\\
        $v_0$ & Systemic velocity & $\mathcal{U}$(-30, 30) & km s$^{-1}$\\
        $\vcirc$ & Max. rotation velocity & $\mathcal{U}$(60, 540) & km s$^{-1}$\\
        $R_\mathrm{vscale}$ & Velocity scale radius & $\mathcal{U}$(0.1, 5.0) & arcsec\\
        $R_h$ & Half-light radius & $\mathcal{U}$(0.1, 5.0) & arcsec\\
        $\log_{10}F_{\mathrm{H}\alpha}$ & H$\alpha$ flux & $\mathcal{U}$(-17, -14) & \shortstack[l]{erg s$^{-1}$\\ cm$^{-2}$}\\
        \tableline
        $q_z$ & Disk aspect ratio & 0.1 & 1 \\
        $\sigma_\lambda^\mathrm{int}$ & Rest-frame line width & 0.05 & nm \\
        $z$ & Redshift & 0.3 & 1 \\
        \tableline
        \end{tabular}
    \caption{KL model parameters. The top section lists the nine parameters inferred by our model, along with the priors from which the training and validation sets are drawn (Sec.~\ref{ssec:datasets}). $\mathcal{U}$(min, max) denotes a uniform distribution. $F_{\mathrm{H}\alpha}$ is the H$\alpha$ flux within the central fiber. The bottom section lists fixed parameters.}
    \label{tab:kl-params}
\end{table}

For our simulations, we use the KL forward-modeling pipeline of \citep{KLInferenceI}, which was modified in \citep{KLDESI} for DESI and Legacy Survey observations. Given a set of galaxy parameters (Table~\ref{tab:kl-params}), the simulator outputs a noiseless $r$-band image and five fiber spectra. We use \textsc{GalSim}~\citep{galsim} to generate an inclined exponential profile for the galaxy's intensity, specified by its half-light radius $R_h$, inclination $i$, and position angle $\theta_\mathrm{int}$, with a fixed aspect ratio of $q_z=0.1$. The profile is then sheared by $(g_1, g_2)$. For the velocity field, we follow \citep{KLDESI} and assume an infinitely thin circular disk with an arctan rotation curve \citep{Courteau97},
\begin{equation}
\label{eqn:vrot}
    v_\mathrm{rot}(R) = \frac{2\vcirc}{\pi}\tan^{-1}\left(\frac{R}{R_\mathrm{vscale}}\right),
\end{equation}
where $R$ is the distance from the center of the disk and $R_\mathrm{vscale}$ is the velocity scale radius. The LoS velocity field is then calculated from $v_\mathrm{rot}$ using $i$ and $\theta_\mathrm{int}$, shifted by the systemic velocity $v_0$, and sheared in the same way as the intensity profile (Eq.~\ref{eqn:shear-vfield}).

For the spectral energy distribution (SED), we use a template galaxy continuum, scaled so that the galaxy's $r$-band AB magnitude is $m_r$, and add a single Gaussian H$\alpha$ emission line on top. \XW{Ask Jiachuan what the continuum template (\texttt{GSB2.spec} in kl-tools) is, and what to cite for it.} As in \citep{KLDESI}, we fix the rest-frame line width to 0.05 nm and place all galaxies at a redshift of $z=0.3$, which puts H$\alpha$ in the $z$ band of the DESI spectrograph. We only simulate a single redshift because we assume that the spectra of an observed galaxy would be cut out in a window centered on its redshifted H$\alpha$ line, so that the position of the line in the spectra does not depend on redshift. The other effects of redshift on our data, namely the apparent size and brightness of the galaxy, are then covered by the ranges of $R_h$, $m_r$, and $F_{\mathrm{H}\alpha}$ in our training set. \XW{Not fully justified until we test the model on samples at other redshifts. Also I think the velocity curve scale per pixel changes with redshift, so measuring on different redshift could bias the kinematics.} The H$\alpha$ flux $F_{\mathrm{H}\alpha}$ is left as a free parameter, which we define as the line flux that falls within the central fiber.

Using the intensity profile, velocity field, and SED, the simulator builds a 3D data cube in position and wavelength, from which it generates the image and spectra; we refer the reader to \citep{KLDESI} for a detailed description of this process. Our image is a $48\times48$ $r$-band stamp with the Legacy Survey pixel scale of 0.2637$\arcsec$, using the DECam $r$-band throughput and a PSF with 1.0$\arcsec$ FWHM. Each fiber spectrum covers a $\sim$5 nm window around the H$\alpha$ line at the DESI dispersion of 0.08 nm per pixel, giving us a $5\times64$ array for the five fibers of each galaxy. The instrument properties and observing configurations follow \citep{KLDESI}, and are listed in Table~\ref{tab:obs-conf}. We generate all images and spectra without noise, and add noise to them during training and inference (Sec.~\ref{ssec:datasets}).

\begin{table}
    \centering
    \begin{tabular}{lccl}
        \toprule
        Property & Image & Fiber & Units\\
        \tableline
        Telescope & CTIO Blanco & KPNO Mayall & - \\
        Effective diameter & 378.29 & 332.42 & cm \\
        Pixel scale & 0.2637 & - & arcsec\\
        PSF FWHM & 1.0 & 1.0 & arcsec \\
        Gain & 4 & 1 & e$^{-}$/ADU \\
        Dispersion & - & 0.08 & nm/pixel \\
        Exposure time & 60 & 180/600 & s \\
        Fiber radius & - & 0.75 & arcsec \\
        Fiber offset & - & 1.5 & arcsec \\
        Data dimensions & $48\times48$ & $5\times64$ & pixel \\
        \tableline
        \end{tabular}
    \caption{Instrument properties and observing configurations used to simulate our images and spectra. Fields with ``-'' are not applicable. Fiber exposure times are listed for the central/offset fibers respectively.}
    \label{tab:obs-conf}
\end{table}

For each galaxy, we place one fiber on its center, two on either side of its observed major axis, and two on either side of its observed minor axis, all with a 1.5$\arcsec$ offset from the center. Following \citep{KLDESI}, the central fiber has an exposure time of 180 s, and each offset fiber has an exposure time of 600 s. Since fiber placements are defined on the observed galaxy, we need to find the observed axes of the galaxy from its intrinsic orientation and shear. We do so by defining the matrix
\begin{equation}
\label{eqn:fiber-transform}
    \mathbf{T} = \mathbf{G}\mathbf{R}\mathbf{P},
\end{equation}
with
\begin{equation}
\label{eqn:fiber-matrices}
    \begin{aligned}
    \mathbf{G} &= \begin{pmatrix}
    1+g_1 & g_2 \\
    g_2 & 1-g_1
    \end{pmatrix}, \\
    \mathbf{R} &= \begin{pmatrix}
    \cos\theta_{int} & -\sin\theta_{int} \\
    \sin\theta_{int} & \cos\theta_{int}
    \end{pmatrix}, \\
    \mathbf{P} &= \begin{pmatrix}
    1 & 0 \\
    0 & \cos i
    \end{pmatrix}.
    \end{aligned}
\end{equation}
Here, $\mathbf{P}$ accounts for the inclination of the disk, $\mathbf{R}$ for its position angle, and $\mathbf{G}$ for the shear. Note that $\mathbf{G}$ is proportional to the inverse of $\mathbf{A}$ in Eq.~(\ref{eqn:shear-mat}), since $\mathbf{T}$ goes from the intrinsic galaxy to the observed image. We then perform a singular value decomposition (SVD) on $\mathbf{T}$, and use its left singular vectors as the observed major and minor axes, onto which we rotate our fiber offsets.

Since multiple solutions exist for the SVD algorithm due to the sign being arbitrary, we need two more conditions to fix them. First, the major axis should point in the same direction as the intrinsic major axis after the transformation, i.e. it should have a positive projection on $\mathbf{T}\,[1, 0]^T$. Second, the minor axis should be $90\degree$ counterclockwise from the major axis. This way, the order of the fibers is tied to the intrinsic orientation of the galaxy, instead of the arbitrary signs returned by the SVD.

\subsection{Training and Validation Data}
\label{ssec:datasets}

We generate 1 million galaxies for our training set and 100 thousand galaxies for our validation set. \XW{Current results are from a model trained on 100k galaxies and validated on 10k. Will retrain on 1M/100k.} For both sets, we sample the nine parameters in Table~\ref{tab:kl-params} from uniform distributions, using Latin Hypercube sampling to achieve even sampling across parameter space. We sample the inclination in $\cos i$ instead of $\sin i$, since $\cos i$ is uniformly distributed for randomly oriented disks, and we sample the H$\alpha$ flux uniformly in $\log_{10}F_{\mathrm{H}\alpha}$. The $r$-band magnitude $m_r$, which sets the brightness of the continuum, is sampled separately from $\mathcal{U}(15, 23.4)$, where the upper limit corresponds to the median $5\sigma$ depth of the Legacy Survey $r$-band images for an extended source with a 0.45$\arcsec$ half-light radius \citep{KLDESI}. The parameter ranges are chosen to be wide enough to contain the galaxies in our test sets (Sec.~\ref{ssec:test-set}). The data generation pipeline is highly parallelizable, using a total of $\sim$2000 CPU hours for the training set.

As mentioned in Sec.~\ref{ssec:kl-model}, we store our images and spectra in their noiseless form, and add noise to them on-the-fly during training and inference. This lets us draw new noise levels and noise realizations for every galaxy at every epoch, so that the model learns from a variety of noise levels for any set of parameters. It also keeps our model as agnostic to observing conditions as possible, since the noise level is specified directly by the signal-to-noise ratio (SNR) rather than by exposure time or sky level. We use a Gaussian noise model for computational efficiency, which we believe is an adequate approximation for ground-based surveys for the current scope of this work.

The level of noise is set by two SNRs: the image SNR $\rho_I$ and the central-fiber H$\alpha$ SNR $\rho_S$. We define $\rho_I$ as the SNR of the image as a whole, where each pixel contributes according to how much of the galaxy's signal it contains:
\begin{equation}
\label{eqn:snr-image}
    \rho_I = \frac{\lVert \bm{M}_\mathrm{I} \rVert}{\sigma_I},
\end{equation}
where $\lVert \bm{M}_\mathrm{I} \rVert$ is the L2 norm of the noiseless image and $\sigma_I$ is the standard deviation of the noise in each pixel. Similarly, we define $\rho_S$ on the H$\alpha$ line in the central fiber,
\begin{equation}
\label{eqn:snr-spec}
    \rho_S = \frac{\lVert \bm{M}_{\mathrm{S},c} - \tilde{M}_{\mathrm{S},c} \rVert}{\sigma_S},
\end{equation}
where $\bm{M}_{\mathrm{S},c}$ is the noiseless central fiber spectrum, $\tilde{M}_{\mathrm{S},c}$ is its median, which we use as an estimate of the continuum level, and $\sigma_S$ is the noise standard deviation in each pixel of the central fiber. Since the noise in counts grows with the square root of exposure time, the noise in each offset fiber is scaled up by $\sqrt{600/180}$ relative to the central fiber. During training, we draw $\rho_I \sim \mathcal{U}(10, 1000)$ and $\rho_S \sim \mathcal{U}(1, 150)$ independently for each galaxy at every epoch. For the validation set, we keep the SNRs and noise realizations fixed across epochs, so that the validation loss can be compared between epochs.

Besides the images and spectra, we also pass a set of metadata into our model: the positions of the five fibers on the image, $m_r$, $\rho_I$, and $\rho_S$. \XW{We pass the true $m_r$ and SNRs as metadata. In real data these would be measured with some uncertainty.} Before passing our data into the model, we normalize each image by its L2 norm, and the five spectra of each galaxy by their combined L2 norm, which keeps the relative flux between fibers. The fiber positions are divided by the fiber offset of 1.5$\arcsec$, and $m_r$, $\rho_I$, and $\rho_S$ are linearly rescaled to $[-1, 1]$ using their training ranges. The purpose of this is to stabilize training and keep our model robust against different observing configurations. As an example, different exposure times on the same galaxy would give different absolute pixel fluxes in the image, but should have no effect on the final shear measurement. Instead, the brightness of the galaxy and the quality of the data are available to the model through $m_r$ and the SNRs.

\subsection{DESI-like Test Sets}
\label{ssec:test-set}

In our final diagnostics, we would like to test our model on a galaxy population that matches as much as possible the one that would be observed by a DESI-KL survey. To best compare with previous results, we build our test sets from the galaxy samples selected in \citep{KLDESI}, specifically samples 3 and 5. Both samples are selected from the DESI Bright Galaxy Survey (BGS) \citep{BGS}, with sample 3 containing all BGS galaxies (BGS Any) and sample 5 containing only BGS Bright galaxies. Both samples also require a disk-like morphology in the Legacy Survey imaging \citep{LSImaging} and a half-light radius of $R_h \geq 1\arcsec$. Sample 5 is the more stringent of the two, giving brighter galaxies but a lower number density. We will refer to the test sets built from these samples as Xu3 and Xu5.

We do not use sample 1 of \citep{KLDESI}, which selects Legacy Survey galaxies with $r \leq 22$. Our test sets require an H$\alpha$ SNR for each galaxy, which is only available for galaxies with a DESI spectrum, and most sample 1 galaxies do not have one. After requiring an H$\alpha$ measurement, the remaining sample 1 galaxies are very similar to those in sample 3, so we do not expect sample 1 to give us additional information. \XW{After all cuts below, Xu1 has 2.42 million eligible galaxies out of 36.5 million catalog galaxies, versus 2.32 million out of 10.4 million for Xu3. Is there a way to get around this for Xu1?} We also do not use samples 2, 4, and 6, which apply a stricter cut of $R_h \geq 1.5\arcsec$, as we did not see significant performance changes in our model between these two cuts.

For each sample, we compile a catalog of all galaxies that pass its selection cuts, and keep only the ones with an H$\alpha$ flux measurement from \textsc{FastSpecFit} \citep{FastSpecFit} in DESI Data Release 1 \citep{DESIDR1}. This catalog gives us empirical joint distributions of $m_r$, $R_h$, image SNR, and H$\alpha$ SNR (see Fig.~\ref{fig:test-set-samples}). We take the image SNR to be the SNR of the Legacy Survey $r$-band flux, and the H$\alpha$ SNR to be that of the \textsc{FastSpecFit} H$\alpha$ flux, and use them as our $\rho_I$ and $\rho_S$ respectively. We further restrict the catalog to the ranges covered by our training set, $15 \leq m_r \leq 23.4$, $0.1\arcsec \leq R_h \leq 5\arcsec$, $10 \leq \rho_I \leq 1000$, and $1 \leq \rho_S \leq 150$, which leaves 2.3 million galaxies in sample 3 and 1.7 million galaxies in sample 5.

We then randomly draw 100 thousand galaxies from each catalog, taking $m_r$, $R_h$, $\rho_I$, and $\rho_S$ from the same galaxy so that the correlations between them are preserved. For each galaxy, we sample $\vcirc$ from the TFR given $m_r$ (Sec.~\ref{ssec:tf-prior}), so that the rotation velocities in our test sets follow a realistic relation with brightness. The rest of the KL parameters, $(g_1, g_2, \theta_\mathrm{int}, \cos i, v_0, R_\mathrm{vscale}, F_{\mathrm{H}\alpha})$, are drawn from the same priors as our training set (Table~\ref{tab:kl-params}). As with our training set, all test galaxies are simulated at $z=0.3$ with the observing configuration of Sec.~\ref{ssec:kl-model}, regardless of their catalog redshift. The resulting test sets have median values of $m_r=19.2$, $R_h=1.6\arcsec$, $\rho_I=174$, and $\rho_S=26$ for Xu3, and $m_r=19.0$, $R_h=1.8\arcsec$, $\rho_I=204$, and $\rho_S=27$ for Xu5.

\begin{figure*}
    \includegraphics[width=\linewidth]{figs/xu_samples_3_5_distributions.png}
    \caption{Redshift, $r$-band magnitude, half-light radius, and H$\alpha$ SNR distributions of samples 3 and 5 from \citep{KLDESI}, for galaxies with H$\alpha$ measurements in DESI DR1. We draw $m_r$, $R_h$, image SNR, and H$\alpha$ SNR jointly from these catalogs to build our Xu3 and Xu5 test sets. The catalog redshift is shown for reference only, as all test galaxies are simulated at $z=0.3$. \XW{Check that this figure is made from the catalogs after the support cuts.}}
    \label{fig:test-set-samples}
\end{figure*}
